\section{Introdução}
\label{sec:introducao}

O avanço dos modelos de linguagem (\textit{Large Language Models}, LLMs) transformou o paradigma de interação entre humanos e máquinas \cite{mcpAtFirstGlance}. No entanto, a ausência inicial de padrões de comunicação entre LLMs e ferramentas externas gerou uma fragmentação técnica que motivou o surgimento do \textit{Model Context Protocol} (MCP), um protocolo aberto criado para padronizar a integração de modelos de IA com programas e fontes de dados externas \cite{shi2026description}. A adoção do protocolo tem crescido de forma acelerada, e um levantamento publicado no início de 2026 apontou que seus Kits de Desenvolvimento de Software (\textit{Software Development Kit}, SDK) já registravam aproximadamente 97 milhões de \textit{downloads} mensais \cite{Avaya2026}. O principal diferencial do MCP reside no uso de ferramentas (\textit{tools}) descritas em linguagem natural, que o agente pode selecionar e invocar em tempo de execução conforme a tarefa a ser resolvida. Esse funcionamento impõe, contudo, uma dependência crítica da qualidade documental, uma vez que o modelo não tem acesso ao código que implementa as ferramentas. As decisões de usá-las, como invocá-las e quais parâmetros fornecer são tomadas com base em suas descrições \cite{shi2026description, 2025arXiv250814704L}. Dessa forma, assim como \textit{code smells} comprometem a manutenibilidade de \textit{softwares} tradicionais, descrições ambíguas ou incompletas comprometem o comportamento do agente, podendo induzir erros de planejamento, uso incorreto de parâmetros e alucinações \cite{2026arXiv260214878M}.

Estudos recentes investigam a prevalência de anomalias (ou \textit{smells}) em descrições de ferramentas MCP por meio de análise documental \cite{2026arXiv260214878M}\textcolor{orange}{\cite{guo2025measurement, shi2026description}}. Essa abordagem avalia se a documentação comunica adequadamente o propósito, os parâmetros e as restrições de uso de uma ferramenta, sem considerar a implementação técnica \cite{mcpAtFirstGlance}. Entretanto, falhas relacionadas à configuração e à documentação representam 38,66\% dos problemas observados em servidores MCP \cite{Taraghi2026}. No MCP, a descrição é um artefato redigido manualmente para um código já existente, e é o código que define o comportamento efetivo da ferramenta; a prevalência de inconsistências entre os dois em servidores reais \cite{shi2026description} mostra que eles podem divergir. Diante disso, existe uma lacuna na \textbf{análise da qualidade das descrições das ferramentas MCP, que, sem o contexto do código, pode não ser suficiente para a avaliação de sua qualidade}.

Essa lacuna é relevante porque inconsistências semânticas podem reduzir a eficácia dos agentes e ampliar a superfície de exposição a vulnerabilidades, incluindo ataques de injeção de comandos \cite{Taraghi2026}\textcolor{orange}{, e relatórios de segurança alertam que ambiguidades documentais podem favorecer a execução de ações não autorizadas ou a exposição indevida de dados sensíveis} \cite{NSA2026}.

Tendo em vista essa lacuna, este trabalho tem como objetivo geral \textbf{investigar se a avaliação da qualidade de descrições de ferramentas MCP é influenciada pela disponibilização do código da ferramenta ao avaliador.} Para alcançar esse objetivo, definem-se os seguintes objetivos específicos: \textcolor{blue}{(i) Examinar se o fornecimento do código da \textit{tool} resulta em maior precisão na avaliação da qualidade de sua descrição; (ii) Identificar os atributos de avaliação mais afetados pela adição do código da \textit{tool}; e (iii) Realizar uma análise estatística sobre os resultados obtidos nas avaliações.}

\textcolor{orange}{Esses objetivos são orientados pelas seguintes questões de pesquisa: \textbf{RQ1}: o fornecimento do código da \textit{tool} ao avaliador aumenta a precisão da avaliação da qualidade de sua descrição, em relação à avaliação baseada apenas na descrição? \textbf{RQ2}: quais atributos da rubrica de qualidade documental são mais afetados pela adição do código da \textit{tool}? \textbf{RQ3}: as diferenças observadas entre os escores dos cenários com e sem código são estatisticamente significativas?}

\textcolor{orange}{Para respondê-las, as ferramentas de 206 servidores MCP de código aberto são avaliadas por três modelos de linguagem atuando como avaliadores (\textit{LLM-as-a-Judge}), que pontuam, em escala Likert de cinco pontos, as seis dimensões da rubrica de qualidade documental de Hasan et al. (2026)\nocite{2026arXiv260214878M}, com e sem acesso ao código da ferramenta.}

As próximas seções deste artigo estão organizadas da seguinte forma. A Seção~\ref{sec:fundamentacao} apresenta a Fundamentação Teórica. A Seção~\ref{sec:trabalhos_relacionados} discute os trabalhos relacionados. A Seção~\ref{sec:materiais} descreve os materiais e métodos\textcolor{orange}{, a Seção~\ref{sec:resultados} apresenta os resultados, a Seção~\ref{sec:ameacas} discute as ameaças à validade e a Seção~\ref{sec:conclusao} apresenta as conclusões e os trabalhos futuros}.
